\section{Inductive types from polynomial endofunctors}
\label{sec:polynomial-types}

This section loosely follows \citet{nunes2023}'s presentation, restricted to inductive types. They take
types over a kinding context $\Delta = \alpha_1 : \mathsf{type}, \ldots, \alpha_n : \mathsf{type}$ so that a
polynomial body $\tau$ in $\mu\alpha.\tau$ may mention other free type variables and a single rule schema
covers every instantiation, whereas our types are closed and each $F \in \Poly_{\cat{C}}$ mentions specific
objects of $\cat{C}$. Where they admit arbitrary type expressions $\mu\alpha.\tau$ and impose strict
positivity as a side-condition on $\tau$, we restrict $\mu$ to elements of $\Poly_{\cat{C}}$, which is
strictly positive by construction. \figref{polynomial-types:syntax} summarises the additional type and term
formers.

\begin{figure}
\[
\begin{array}{rclr}
   P, Q &::=& \const(\sigma) \mid \mathsf{var} \mid P + Q \mid P \times Q & \text{(polynomials over }\cat{C}\text{)} \\
   \tau, \sigma &::=& \cdots \mid \mu P & \text{(types)} \\
   t, s &::=& \cdots \mid \roll\,t \mid \fold\,t\,\with\,x.s & \text{(terms)}
\end{array}
\]
\caption{Additional syntactic forms for polynomial types.}
\label{fig:polynomial-types:syntax}
\end{figure}

\subsection{The category $\Poly_{\cat{C}}$ of polynomial endofunctors}
\label{sec:polynomial-types:poly}

In the style of \citet[Definition~6]{nunes2023} we treat the polynomial endofunctors as a (sub)category of
$\Cat$, generated under a fixed collection of operations.

\begin{definition}[Polynomial endofunctors]
\label{def:polynomial-types:Poly}
Let $\cat{C}$ be a category with finite products $(\times, 1)$ and finite coproducts $(\coprod, 0)$. The
category $\Poly_{\cat{C}}$ is the smallest subcategory of $\Cat$ satisfying:
\begin{itemize}
\item $\cat{C}$ is an object of $\Poly_{\cat{C}}$;
\item for any object $A$ of $\cat{C}$, the constant functor $\const_A: \cat{C} \to \cat{C}$ is a morphism of
$\Poly_{\cat{C}}$;
\item the identity functor $\cat{C} \to \cat{C}$ is a morphism of $\Poly_{\cat{C}}$;
\item if $F, G: \cat{C} \to \cat{C}$ are morphisms of $\Poly_{\cat{C}}$ then so are the pointwise coproduct
$F \coprod G$ and pointwise product $F \times G$.
\end{itemize}
We say $\cat{C}$ \emph{has $\Poly$-types} if every endofunctor in $\Poly_{\cat{C}}$ admits an initial
algebra.
\end{definition}

\noindent The endofunctors of $\Poly_{\cat{C}}$ are thus exactly those generated by constants, the identity,
and the pointwise binary product and coproduct. Unlike \citet{nunes2023}, we do not close $\Poly_{\cat{C}}$
under $\mu$, so $\mu F$ is an object of $\cat{C}$ rather than a polynomial; this loses no expressivity,
since $\mu F$ can subsequently appear as a $\const$ in further polynomials.

\subsection{Typing rules}
\label{sec:polynomial-types:language}

Here $P(\tau)$ denotes the type obtained by substituting $\tau$ for $\mathsf{var}$ in $P$, with
$\const(\sigma)(\tau) = \sigma$, $\mathsf{var}(\tau) = \tau$, $(P + Q)(\tau) = P(\tau) + Q(\tau)$, and
$(P \times Q)(\tau) = P(\tau) \times Q(\tau)$. The typing rules for the new term formers are
\[
\begin{array}{c}
   \infer{\Gamma \vdash \roll\,t : \mu P}{\Gamma \vdash t : P(\mu P)} \qquad
   \infer{\Gamma \vdash \fold\,t\,\with\,x.s : \tau}{\Gamma \vdash t : \mu P \qquad
   \Gamma, x : P(\tau) \vdash s : \tau}
\end{array}
\]
The second rule uses the parametric (open) form of the catamorphism: the algebra body $s$ lives in the
extended context $\Gamma, x : P(\tau)$ rather than as a closed function $P(\tau) \to \tau$, avoiding the
need for exponentials in $\cat{C}$.

\subsection{Existence of initial algebras in $\Fam(\cat{C})$}
\label{sec:polynomial-types:fam}

\citet[Proposition~81]{nunes2023} establish that $\Fam(\Set)$ has initial algebras for their polynomials. We
give a direct construction in $\Fam(\cat{C})$ via Martin-L\"of $W$-types
\citep[Wellorderings, pp.~43--47]{martinLof1984}, restricted to the finitary fragment (where each shape has
finite arity), which suffices because our polynomials are built from finite $+$ and $\times$.

\begin{proposition}
\label{prop:polynomial-types:fam-has-poly}
If $\cat{C}$ has finite products $(\times, 1)$ then $\Fam(\cat{C})$ has $\Poly$-types.
\end{proposition}

\noindent The construction proceeds by induction on the syntactic polynomial $P$ with $\sem{P} = F$, where
$\sigma$ in each $\const(\sigma)$ interprets to a pair $(X, \partial X) \in \Fam(\cat{C})$. The carrier
$\mu F = (W_P, \partial W_P) \in \Fam(\cat{C})$ has:
\begin{itemize}
\item index set $W_P$, the set of $P$-shaped well-founded trees: the tree has the shape of $P$, with each
$\const(\sigma)$ position carrying an element of $X$ and each $\mathsf{var}$ position carrying a recursive
subtree;
\item fibre $\partial W_P(t) \in \cat{C}$ for each tree $t \in W_P$, defined by structural recursion on $P$:
$\partial X_x$ at each $\const(\sigma)$ position labelled $x$; the recursive fibre at each $\mathsf{var}$
position; the chosen branch's fibre at each $+$ position; and the product of branch fibres at each $\times$
position.
\end{itemize}
The algebra map $(\inMap_F, \partial \inMap_F): F(\mu F) \to \mu F$ has $\inMap_F$ assembling a single tree
from its constituent parts (const labels, branch choices, recursive subtrees), and $\partial \inMap_F$ the
identity because the fibre of the assembled tree is computed by the same recursion as $F$'s action on the
parts.

Initiality (proved by induction on the tree) unpacks as the standard $\beta$ and $\eta$ equations. For any
$F$-algebra $\alpha: F(Y) \to Y$, write $h_\alpha: \mu F \to Y$ for the unique algebra morphism. Then
\begin{align*}
h_\alpha \circ \inMap_F &= \alpha \circ F(h_\alpha) \tag{$\beta$} \\
h_{\inMap_F} &= \id_{\mu F} \tag{$\eta$}
\end{align*}
where $\beta$ is the algebra-morphism property of $h_\alpha$ and $\eta$ applies uniqueness with $\alpha =
\inMap_F$, where $\id_{\mu F}$ is trivially an algebra morphism. At the term level, $\beta$ reduces
$\fold\,(\roll\,t)\,\with\,x.s$ by recursively folding each $\mu P$-child of $t$ before substituting into
$s$, and $\eta$ says $\fold\,t\,\with\,x.\roll\,x = t$.

